\documentclass[10pt]{article}

\usepackage[margin=0.65in]{geometry}
\usepackage{amsmath,amssymb}
\usepackage{graphicx}
\newcommand{\R}{\mathbb{R}}
\newcommand{\down}{\ensuremath{\downarrow}}
\newcommand{\up}{\ensuremath{\uparrow}}
\newcommand{\asset}[1]{\includegraphics[width=\linewidth]{results_assets/#1}}

\title{DriftXpress MNIST Results}
\author{}
\date{}

\begin{document}
\maketitle

\section*{Experimental record}

This report was assembled from the saved code, textual outputs, and embedded
figures in the five notebooks in this folder.  No notebook was executed or
modified while producing this report.  The experiment uses MNIST digits
\(\{0,8,2,4,7,9\}\), with 35,788 training examples, 6,005 test examples,
batch size 500, latent dimension 32, 200 drift landmarks, and 100 configured
training epochs.  Because the notebook loops use \texttt{range(num\_epochs+1)},
the saved runs actually include epochs 0 through 100.

\paragraph{Metric comparability.}
All five evaluation cells compare decoded generated images against flattened
MNIST test pixels.  The reported ``FID'' is the implementation in
\texttt{src/eval.py}, not canonical Inception-FID; PCA and UMAP are visual
projections in these notebooks and do not produce a scalar score.

\section*{Aggregate evaluation metrics}

\begin{table}[htbp]
\centering
\small
\caption{Standardized pixel-space evaluation. Lower is better for SW, MMD, and FID; higher is better for coverage and entropy.}
\begin{tabular}{l l rrrrr}
\hline
Algorithm & evaluation space & SW $\down$ & MMD $\down$ & FID $\down$ & coverage $\up$ & entropy $\up$ \\
\hline
1 & pixels & 0.156390 & 0.142476 & 32.506738 & 1.000000 & 0.999951 \\
2 & pixels & 0.055698 & 0.019198 & 9.749420 & 1.000000 & 0.999930 \\
3 & pixels & 0.145378 & 0.150643 & 60.209453 & 1.000000 & 0.999966 \\
4 & pixels & 0.088063 & 0.043731 & 15.564459 & 1.000000 & 0.999738 \\
5 & pixels & 0.115568 & 0.095509 & 49.368244 & 1.000000 & 0.999588 \\
\hline
\end{tabular}
\end{table}

\begin{table}[htbp]
\centering
\scriptsize
\caption{Per-digit pixel-space FID values.}
\begin{tabular}{c crrrrr}
\hline
Digit & space & Alg. 1 & Alg. 2 & Alg. 3 & Alg. 4 & Alg. 5 \\
\hline
0 & pixel & 58.9163 & 13.2347 & 86.5164 & 29.0469 & 82.2101 \\
8 & pixel & 52.1156 & 17.3052 & 80.4397 & 24.8819 & 67.9792 \\
2 & pixel & 46.5633 & 19.7732 & 74.4471 & 25.1308 & 66.0384 \\
4 & pixel & 39.0935 & 16.1500 & 66.0835 & 20.6945 & 52.5338 \\
7 & pixel & 34.4265 & 11.9323 & 68.3462 & 9.6142 & 47.2885 \\
9 & pixel & 39.9167 & 14.6916 & 69.6437 & 18.5663 & 51.2729 \\
\hline
\end{tabular}
\end{table}

\section*{Algorithm summaries and losses}

\subsection*{Algorithm 1: joint conditional training}

Algorithm 1 trains the VAE and conditional latent generator in one loop with
one optimizer object.  Its drift term acts directly on the raw generator
latent, so the drift gradient reaches the generator but bypasses the encoder
and decoder.  The saved run has pixel FID 32.51 and perfect class coverage,
although its cache is built before training from sampled encoder latents.

\[
\mathcal L_1=\mathcal L_{\mathrm{recon}}
 +\lambda_{\mathrm{KL}}\mathcal L_{\mathrm{KL}}
 +\lambda_{\mathrm d}\mathcal L_{\mathrm{drift}}(z^-)
 +\lambda_{\mathrm v}\bigl(\mathcal L_{\mathrm{var}}(z^+)+\mathcal L_{\mathrm{var}}(z^-)\bigr)
 +\lambda_{\mathrm c}\bigl(\mathcal L_{\mathrm{cov}}(z^+)+\mathcal L_{\mathrm{cov}}(z^-)\bigr).
\]

\subsection*{Algorithm 2: separate VAE and generator optimization}

Algorithm 2 first trains the VAE and then freezes it while optimizing a
separate conditional latent generator.  The generator is trained only with
drift, variance, and covariance objectives, so it cannot change the learned
decoder or encoder representation.  Its pixel-space evaluation gives FID
9.75 with complete class coverage.

\[
\begin{aligned}
\mathcal L_{\mathrm{VAE}}&=\mathcal L_{\mathrm{recon}}
 +\lambda_{\mathrm{KL}}\mathcal L_{\mathrm{KL}}
 +\lambda_{\mathrm v}\mathcal L_{\mathrm{var}}(\mu)
 +\lambda_{\mathrm c}\mathcal L_{\mathrm{cov}}(\mu),\\
\mathcal L_{\mathrm G}&=\lambda_{\mathrm d}\mathcal L_{\mathrm{drift}}(z^-)
 +\lambda_{\mathrm v}\mathcal L_{\mathrm{var}}(z^-)
 +\lambda_{\mathrm c}\mathcal L_{\mathrm{cov}}(z^-).
\end{aligned}
\]

\subsection*{Algorithm 3: fully joint re-encoded drift}

Algorithm 3 sends generated latents through the decoder and deterministic
encoder mean before applying drift.  This creates a differentiable path from
the re-encoded drift loss to the generator, decoder, and encoder, without a
cycle loss.  The saved pixel FID is 60.21, and the generated digit grid shows
that the re-encoding path alone does not sufficiently constrain the raw
generator latent.

\[
\mathcal L_3=\mathcal L_{\mathrm{recon}}
 +\lambda_{\mathrm{KL}}\mathcal L_{\mathrm{KL}}
 +\lambda_{\mathrm d}\mathcal L_{\mathrm{drift}}(\widetilde z^-)
 +\lambda_{\mathrm v}\bigl(\mathcal L_{\mathrm{var}}(\mu)+\mathcal L_{\mathrm{var}}(\widetilde z^- )\bigr)
 +\lambda_{\mathrm c}\bigl(\mathcal L_{\mathrm{cov}}(\mu)+\mathcal L_{\mathrm{cov}}(\widetilde z^- )\bigr),
\qquad
\widetilde z^-=\mu_\phi(D_\psi(G_\theta(\varepsilon,c))).
\]

\subsection*{Algorithm 4: latent discriminator}

Algorithm 4 adds a class-conditional latent discriminator to Algorithm 1.
The discriminator is trained with detached real and generated latents, while
the generator receives a non-saturating BCE adversarial loss through the
temporarily frozen discriminator.  Its pixel-space evaluation gives FID
15.56 with complete class coverage.

\[
\begin{aligned}
\mathcal L_4&=\mathcal L_1+\lambda_{\mathrm{adv}}\mathcal L_{\mathrm{adv}},\\
\mathcal L_{\mathrm{disc}}&=\tfrac12\left[
 \operatorname{BCE}(D_\omega(\operatorname{sg}(z^+),c),1)
 +\operatorname{BCE}(D_\omega(\operatorname{sg}(z^-),c),0)\right],\\
\mathcal L_{\mathrm{adv}}&=\operatorname{BCE}(D_\omega(z^-,c),1).
\end{aligned}
\]

\subsection*{Algorithm 5: pretraining plus raw and re-encoded drift}

Algorithm 5 pretrains the VAE for 25 epochs and builds the cache from
deterministic encoder means.  During fully joint training it averages drift
on the raw generator latent and on the decoded--re-encoded latent, preserving
both direct generator control and the connected gradient path.  The L1
coefficient is present but set to zero, and the saved pixel metrics are
SW 0.115568, MMD 0.095509, and FID 49.37.

\[
\begin{aligned}
\mathcal L_{\mathrm{pre}}&=\mathcal L_{\mathrm{recon}}
 +\lambda_{\mathrm{L1}}\mathcal L_1
 +\lambda_{\mathrm{KL}}\mathcal L_{\mathrm{KL}},\\
\mathcal L_5&=\mathcal L_{\mathrm{recon}}
 +\lambda_{\mathrm{L1}}\mathcal L_1
 +\lambda_{\mathrm{KL}}\mathcal L_{\mathrm{KL}}
 +\lambda_{\mathrm d}\tfrac12\left[
 \mathcal L_{\mathrm{drift}}(z^-)+\mathcal L_{\mathrm{drift}}(\widetilde z^-)
 \right]\\
&\quad+\lambda_{\mathrm v}\bigl(\mathcal L_{\mathrm{var}}(\mu)+\mathcal L_{\mathrm{var}}(\widetilde z^- )\bigr)
 +\lambda_{\mathrm c}\bigl(\mathcal L_{\mathrm{cov}}(\mu)+\mathcal L_{\mathrm{cov}}(\widetilde z^- )\bigr).
\end{aligned}
\]

\section*{Model parameters and optimization settings}

The notebook parameter-count cells instantiate constructor defaults rather
than the later configured \(\texttt{LATENT\_DIM}=32\).  The table below gives
the parameter counts for the actual trained configuration, recomputed from
the architectures used by each training loop.

\begin{table}[htbp]
\centering
\scriptsize
\caption{Architecture and configured optimization parameters.}
\begin{tabular}{p{0.06\textwidth} p{0.15\textwidth} p{0.29\textwidth} p{0.32\textwidth} p{0.13\textwidth}}
\hline
Algorithm & trainable parameters & architecture & optimization settings & activations \\
\hline
1 & 357,280 & $E$: 784--192--96--32; $D$: 32--96--192--784; $G$: 48--48--48--32; label embedding 16 & $T=2$, $r=200$, $\lambda_{KL}=10^{-5}$, $\lambda_d=30$, $\lambda_v=.1$, $\lambda_c=.3$, AdamW $10^{-3}$ & LeakyReLU encoder/decoder, SiLU generator, Sigmoid output \\
2 & 357,280 total; VAE 348,560, $G$ 8,720 & Same VAE and generator as Algorithm 1; VAE frozen in Stage II & $T=2$, $r=200$, same loss weights; 100 VAE epochs then 100 generator epochs, AdamW $10^{-3}$ & LeakyReLU, SiLU, Sigmoid \\
3 & 357,280 & Same VAE and generator; generated path includes $G\to D\to\mu$ & $T=2$, $r=200$, same loss weights; AdamW $10^{-3}$ & LeakyReLU, SiLU, Sigmoid \\
4 & 357,280 main model + 7,457 discriminator & Algorithm 1 model plus $D_\omega$: 48--64--64--1 with label embedding 16 & $T=2$, $r=200$, $\lambda_{adv}=1$, model/discriminator LR $10^{-3}$ & LeakyReLU discriminator; BCE-with-logits \\
5 & 357,280 & Same VAE and generator; deterministic $\mu$ cache; raw and re-encoded drift & 25 VAE pretraining epochs, then 100 joint epochs; $T=2$, $r=200$, $\lambda_{L1}=0$, AdamW $10^{-3}$ & LeakyReLU, SiLU, Sigmoid \\
\hline
\end{tabular}
\end{table}

\section*{Saved generated MNIST outputs}

Each figure reproduces the embedded notebook output.  The rows are labelled
by the notebook as real images, reconstructions, and generated samples when
all three rows are present.

\begin{figure}[htbp]\centering\asset{algorithm1_digits.png}
\caption{Algorithm 1 saved MNIST output.}\end{figure}
\begin{figure}[htbp]\centering\asset{algo2_digits.png}
\caption{Algorithm 2 saved MNIST output.}\end{figure}
\begin{figure}[htbp]\centering\asset{algo3_digits.png}
\caption{Algorithm 3 saved MNIST output.}\end{figure}
\begin{figure}[htbp]\centering\asset{algo4_digits.png}
\caption{Algorithm 4 saved MNIST output.}\end{figure}
\begin{figure}[htbp]\centering\asset{algo5_digits.png}
\caption{Algorithm 5 saved MNIST output.}\end{figure}

\section*{PCA and UMAP projections}

The first projection for each algorithm uses digit labels; the second uses a
binary real-versus-generated label.  These are visual diagnostics rather than
numeric scores.  The saved Algorithm 1 figures show an empty UMAP panel,
consistent with the previously reported UMAP environment problem; the other
notebooks contain rendered PCA/UMAP panels.

\begin{figure}[htbp]\centering\asset{algorithm1_projection_1.png}
\caption{Algorithm 1 PCA/UMAP projection colored by digit.}\end{figure}
\begin{figure}[htbp]\centering\asset{algorithm1_projection_2.png}
\caption{Algorithm 1 PCA/UMAP projection colored by real/generated source.}\end{figure}
\begin{figure}[htbp]\centering\asset{algo2_projection_1.png}
\caption{Algorithm 2 PCA/UMAP projection colored by digit.}\end{figure}
\begin{figure}[htbp]\centering\asset{algo2_projection_2.png}
\caption{Algorithm 2 PCA/UMAP projection colored by real/generated source.}\end{figure}
\begin{figure}[htbp]\centering\asset{algo3_projection_1.png}
\caption{Algorithm 3 PCA/UMAP projection colored by digit.}\end{figure}
\begin{figure}[htbp]\centering\asset{algo3_projection_2.png}
\caption{Algorithm 3 PCA/UMAP projection colored by real/generated source.}\end{figure}
\begin{figure}[htbp]\centering\asset{algo4_projection_1.png}
\caption{Algorithm 4 PCA/UMAP projection colored by digit.}\end{figure}
\begin{figure}[htbp]\centering\asset{algo4_projection_2.png}
\caption{Algorithm 4 PCA/UMAP projection colored by real/generated source.}\end{figure}
\begin{figure}[htbp]\centering\asset{algo5_projection_1.png}
\caption{Algorithm 5 PCA/UMAP projection colored by digit.}\end{figure}
\begin{figure}[htbp]\centering\asset{algo5_projection_2.png}
\caption{Algorithm 5 PCA/UMAP projection colored by real/generated source.}\end{figure}

\section*{Gradient-flow diagrams}

The current \texttt{formulations.tex} contains the combined gradient-flow
diagram for Algorithms 1--3.  It is included below as page 3 of its compiled
PDF.  The source file does not currently contain separate Algorithm 4 and
Algorithm 5 diagrams, so their flow diagrams are included immediately after
it using the computation graphs implemented in their notebooks.

\begin{figure}[htbp]
\centering
\includegraphics[page=3,width=0.98\textwidth]{formulations.pdf}
\caption{Gradient-flow diagram for Algorithms 1--3, included from the compiled \texttt{formulations.tex}.}
\end{figure}

\subsection*{Algorithm 4 flow}

\begin{figure}[htbp]
\centering
\[
\begin{array}{ccccc}
(x,c)&\xrightarrow{q_\phi}&z^+&\xrightarrow{D_\psi}&\mathcal L_{\rm recon}+\lambda_{\rm KL}\mathcal L_{\rm KL}\\[2mm]
(\varepsilon,c)&\xrightarrow{G_\theta}&z^-&\xrightarrow{\ \mathcal L_{\rm drift,var,cov}\ }&\boxed{\nabla_{\phi,\psi,\theta}\ \text{update}}\\[2mm]
\multicolumn{2}{c}{\operatorname{sg}(z^+),\operatorname{sg}(z^-)}&\xrightarrow{D_\omega}&\mathcal L_{\rm disc}\ \text{(detached)}&\\[2mm]
z^-&\xrightarrow{D_\omega\ \text{with frozen }\omega}&\mathcal L_{\rm adv}&\xrightarrow{\ \text{gradient to }\theta\ }&\boxed{\text{update}}
\end{array}
\]
\caption{Algorithm 4 adds a detached discriminator-training path and an adversarial gradient path to the generator.}
\end{figure}

\subsection*{Algorithm 5 flow}

\begin{figure}[htbp]
\centering
\[
\begin{array}{ccccc}
x&\xrightarrow{q_\phi}&\mu&\xrightarrow{D_\psi}&\mathcal L_{\rm recon}+\lambda_{\rm KL}\mathcal L_{\rm KL}\quad\text{(25-epoch pretraining)}\\[2mm]
\mu&\xrightarrow{\text{deterministic}}&\text{fixed class cache}&&\\[2mm]
(x,c)&\xrightarrow{q_\phi}&z^+&\xrightarrow{D_\psi}&\mathcal L_{\rm recon}+\lambda_{\rm KL}\mathcal L_{\rm KL}\\[2mm]
(\varepsilon,c)&\xrightarrow{G_\theta}&z_{\rm raw}&\xrightarrow{D_\psi,\ \mu_\phi}&\widetilde z^-&\\[2mm]
z_{\rm raw},\widetilde z^-&\xrightarrow{\ \frac12(\mathcal L_{\rm drift}^{\rm raw}+\mathcal L_{\rm drift}^{\rm reenc})\ }&&&\boxed{\nabla_{\phi,\psi,\theta}\ \text{fully joint update}}
\end{array}
\]
\caption{Algorithm 5 first stabilizes the VAE and deterministic cache, then combines raw-latent and re-encoded drift during fully joint training.}
\end{figure}

\end{document}
